\section{INTRODUCTION}
Media engagement is more than entertainment for blind and visually impaired (BVI) people, it’s a means of participation, self-representation, and advocacy \cite{emara_its_2025,lyu_because_2024,niu_please_2024}. Research demonstrates that video creation offers BVI individuals significant wellbeing benefits like enhanced social connection, confidence building, and community identity formation, It also enable them to share their stories, educate audiences, and challenge stereotypes despite platform accessibility barriers \cite{seo_understanding_2021,zhang_understanding_2023}.  Yet historically, BVI communities have been underrepresented and sometimes excluded from mainstream media production and cultural venues  due to inaccessible tools, lack of training opportunities, and systemic ableism. \cite{lyu_because_2024,raymond_examination_2024,lyu_i_2024-1}. Unlike text-based blogging or audio-only podcasting, video has emerged as the dominant medium of contemporary digital communication, fundamentally reshaping how we participate in social and cultural life \cite{ha_global_2022}). At the meantime, video creation presents unique challenges and opportunities for BVI creators including inaccessible editing interfaces, inability to verify visual content, and platform-specific barriers.  While the visual-centric nature of video might seem inherently exclusionary, BVI creators have developed innovative approaches to video production that challenge assumptions about who can create visual media and how \cite{seo_understanding_2021,zhang_understanding_2023}. These creators navigate complex workflows involving scripting, audio recording, visual composition, and editing, often developing workarounds for inaccessible tool s while producing content that resonates with both sighted and non-sighted audiences \cite{xiao_understanding_2025}. 

HCI (Human-Computer Interaction) has substantially improved media \textit{consumption} for BVI audiences, including audio description ecosystems, search and viewing support, and design spaces that prioritise non-visual access \cite{aydin_towards_2020, natalie_audio_2024, wang_toward_2021}. In terms of media \textit{creation}, a growing body of work now examines creator practices and proposes tools for accessible production \cite{huh_avscript_2023,lee_altcanvas_2024}. However, practical training programs and capacity-building initiatives for accessible video creation remain largely absent from educational curricula and community-based learning environments.  While self-taught BVI creators have pioneered individual workarounds and shared knowledge through online communities \cite{lyu_i_2024-1,seo_understanding_2021}, the complexity of video production, requiring simultaneous mastery of inaccessible tools, visual verification, and platform-specific requirements, creates barriers that isolated learning struggles to overcome \cite{zhang_understanding_2023}. 

This training gap is particularly significant when considering the collaborative nature of media production. While existing accessibility research often frames BVI creators as individual users navigating inaccessible systems, video creation inherently involves collaboration , from the division of labour required for complex production tasks (scripting, filming, editing, publishing) to the interdependence between visual and non-visual expertise in composing shots, verifying visual content, and ensuring accessibility features \cite{huh_avscript_2023,zhang_understanding_2023,xiao_understanding_2025}. Recent HCI scholarship has begun exploring ability-diverse collaboration, where participants with different abilities work together as interdependent contributors rather than in helper-helped dynamics \cite{bennett_interdependence_2018,xiao_systematic_2024}. Yet we lack empirical understanding of how to structure learning environments that foster genuine ability-diverse collaboration in creative contexts, moving beyond accessibility accommodations toward collaborative interdependence.

To address this gap, we conducted seven workshops designed to teach non-professional video creation skills to participants with visual impairments. Through this case study approach and in line with a participatory action research approach, we investigate two critical research questions: 

\begin{itemize}
    \item \textbf{RQ1:} What barriers/challenges affect video creation education in an ability-diverse team?
    \item \textbf{RQ2:} What collaboration patterns emerge in video creation education within ability-diverse teams and how could them inform the accessible video creation education?
\end{itemize}

\textbf{Contribution Statement: }

This work makes three primary contributions. First, we provide an \textbf{empirical account of the challenges and opportunities} that arise when teaching video creation in ability-diverse settings, offering insights that extend beyond accessibility to fundamental questions about pedagogical design in HCI. Second, we identify specific \textbf{patterns of collaboration that emerge when participants with different visual abilities work together on video projects}, revealing how ability diversity can enhance rather than hinder creative outcomes. Third, we present \textbf{design implications for future tools and curricula settings that support ability-diverse video creation}, grounded in our observations of how participants navigated technological barriers, developed workarounds, and leveraged peer support. These findings contribute to the growing discourse on creative accessibility in HCI while offering practical guidance for researchers and practitioners working to democratize media production.


